\chapter{Structural Network Operations}\label{ch09:networks}

\section{The capability-gating philosophy}\label{ch09:sec:philosophy}

The network layer is where the Ohm analogy is most seductive and least
earned. The Quantification Guide presents series addition of resistances,
reciprocal aggregation of parallel resistances, and multiplication of stage
gains as working rules --- arithmetic one may simply perform on constraint
values to obtain system-level conclusions. The operationalization addendum,
read against it, reclassifies the same arithmetic as \emph{preregistered
predictions}: falsifiable hypotheses (the gain-composition and redundancy
hypotheses of the registered claim ledger) whose defining rules must be
dropped if the corresponding experiments fail. And Appendix~I applies the
arithmetic to paths while warning, in the same breath, that co-present
guardrails are not parallel routes and that flow remains local. This is the
validation-status conflict preserved in the source audit, and the candidate's
response to it is the organizing principle of this chapter:

\begin{ruleenv}[Structural-only network operations]\label{ch09:rule:structural}
Series equivalence, parallel equivalence, and gain composition exist in the
candidate only as \emph{unauthenticated structural research calculations}.
Each operation is capability-gated: it runs only inside an admitted
operation context, under a reviewed and current calibration profile, on
typed inputs that satisfy coordinate, unit, topology, region, linearity,
and evidence requirements; it emits a result that binds its complete
derivation; and it carries an authorization state that is, in every
implemented mode, explicitly unauthenticated. No network result asserts
that the arithmetic is empirically valid; it asserts only that the
arithmetic was performed, on these children, under these gates, with this
provenance.
\end{ruleenv}

The word ``capability'' is doing precise work here. The implemented
capability enumeration names the network operations
(\texttt{SERIES\_EQUIVALENCE}, \texttt{PARALLEL\_EQUIVALENCE},
\texttt{GAIN\_COMPOSITION}) alongside the prediction and diagnosis
capabilities, and every operation result embeds an operation binding whose
capability field records which gate was exercised. The binding also records
a trust mode --- \texttt{STRUCTURAL\_RESEARCH\_ONLY} or
\texttt{INTEGRITY\_CONFIGURATION} --- and an authorization state that is
locked to the trust mode by validation: structural mode must carry
\texttt{UNAUTHENTICATED\_STRUCTURAL}, integrity mode must carry
\texttt{UNAUTHENTICATED\_INTEGRITY}, and any other pairing is a validation
error. There is no third mode. The type system does not contain an
authenticated authorization state for these operations, so no configuration
of the candidate, however keyed or attested, can produce a network result
that claims authenticated authority (Section~\ref{ch09:sec:ed25519}).

Why gate so heavily what is, after all, elementary arithmetic? Because the
arithmetic is not where the risk lives. Adding two positive decimals is
harmless; \emph{calling the sum an equivalent resistance} is a modeling
claim with unstated preconditions --- that the two components are actually
in series in the relevant causal topology, that their local linear regimes
are simultaneously applicable, that their fitted parameters share
coordinates and units, that their errors combine as declared, and that no
interaction term is being silently zeroed. The electrical formulas are
theorems in circuit theory because Kirchhoff's laws supply those
preconditions; here, nothing supplies them. The candidate therefore inverts
the Guide's presentation: instead of rules that are valid until questioned,
it implements calculations that are inert until every stated precondition
is demonstrated to the type system, and whose outputs remain, even then,
research structures awaiting the preregistered experiments. This is the
verification-validation separation \cite{sargent} applied to network
arithmetic: the software verifies its gates exhaustively; validation of the
composition laws themselves belongs to the registered hypotheses and has
not occurred.

A second, quieter principle governs the layer: \emph{every network
resistance replays}.
Every network result is factory-only --- direct construction is rejected by
token check --- and every result binds its exact ordered children by digest,
so that a verifier holding the result and the evidence registry can
re-resolve the children, re-run the factory, and demand digest equality
with the original. A network value that cannot be replayed to its exact
bytes does not validate. This converts the network layer from a calculator
into a closed derivation system: there are no free-floating equivalent
resistances anywhere in the candidate, only endpoints of replayable
derivation chains.

\section{Context admission and immutability}\label{ch09:sec:context}

Before any network operation runs, an operation context must be admitted,
and the admission mechanics matter enough to the layer's guarantees to be
spelled out here rather than assumed.

A context binds five things: the calibration profile; an evidence
registry; a capability authority (manifest plus verification machinery);
a tuple of capability attestations, possibly empty, with duplicate
capabilities rejected; and the caller-supplied replay time, with an
optional recorded clock receipt. Admission is not a constructor
formality. At admission the context re-validates every object in the
evidence registry; resolves the authority's manifest reference and
requires byte equality between the registered artifact payload and the
manifest's canonical bytes; asserts the profile current at the replay
time (reviewed, drift-current, inside the validity window); requires
every artifact the profile references --- code, dependency lock, corpus,
corpus rights, both splits, drift check, directions, full-scale sources,
controls --- to be resolvable in the registry; and requires the evidence
behind every presented attestation. Only then does it take digest
snapshots: of the profile, the manifest, each attestation, the replay
time, and the clock receipt.

The snapshots are the immutability mechanism. Every subsequent operation
--- series, parallel, gain, prediction, residual --- begins by
re-validating the context against its snapshots: the types must be
unchanged, the registry must re-validate, and the profile digest,
manifest digest, attestation digests, replay time, and clock receipt
must equal their admitted values, with any divergence a validation
error. A context is therefore not a session that can drift; it is a
frozen frame, and every network result records which frame it was
computed in. Two results with the same binding digests were computed
under identical premises; two results with different profile digests are
formally incomparable, and nothing in the layer will combine them ---
the one-profile requirement of the admission rules is checked against
the context's single profile, so cross-profile aggregation is not a
forbidden option but an unexpressible one.

One further mechanical point completes the picture: all network
arithmetic runs inside a declared exact-decimal arithmetic context, so
that sums, reciprocals, products, and square roots are computed at fixed
high precision with deterministic behavior across platforms, and the
digest-equality replay obligation of
Section~\ref{ch09:sec:philosophy} is meaningful --- bit-identical
replay would be unachievable over platform-dependent binary floating
point. The context's precision and rounding mode are implementation
constants participating in the code and dependency-lock identities; the
governing sources do not restate their numeric values, and recording
them in this text is an open documentation item, since the
digest-equality replay guarantee below depends on them.

\section{Series equivalence}\label{ch09:sec:series}

\subsection{The electrical analogy, bracketed}

In circuit theory, resistances in series add because the same current
traverses each element and voltages sum along the path. Both premises are
physical facts about charge conservation and potential. The candidate's
series rule borrows the \emph{arithmetic} --- the equivalent value is the
sum of component values --- while explicitly bracketing the premises: there
is no conserved current in a steered activation pathway, no potential
whose drops sum, and no theorem that two locally linear response segments
compose into one. The bracketing is not rhetorical. It is implemented as
the requirement set below, which demands precisely the structural
conditions under which the sum is at least \emph{coherent} (same
coordinates, same units, same profile, contiguous topology), while the
question of whether the sum \emph{predicts} anything is left, deliberately
unanswered, to the registered redundancy and composition experiments. Until
those close, a series equivalent is a bookkeeping object: the sum of some
fitted parameters, with provenance.

\subsection{Requirements as implemented}

\begin{ruleenv}[Series admission]\label{ch09:rule:series}
A series equivalence over components $C_1, \ldots, C_n$ ($n \ge 2$) is
admitted only when all of the following hold.
\begin{enumerate}
\item \emph{Typed components.} Each $C_i$ is an exact resistance-component
object: a component identity, a resistance measurement, a source node, a
target node, a declared region, and a declared linearity coefficient.
Duplicate component identities, duplicate component digests, and duplicate
underlying resistance digests are all rejected --- the same fitted
resistance cannot be counted twice under two names.
\item \emph{Calibrated or replayable resistances.} Each component's
resistance is of quantity kind \texttt{EFFECTIVE\_RESISTANCE} with role
\texttt{CALIBRATED\_HELD\_OUT} or \texttt{NETWORK\_EQUIVALENT}; apparent
same-sample and design-proposal roles are inadmissible. A
network-equivalent child is not taken on faith: it is recursively replayed
through its own factory and digest-checked before use.
\item \emph{One unit.} Every resistance is in the single unit
\texttt{ALPHA\_PER\_ACTIVATION}; the equivalent is emitted in the same
unit. No unit conversion exists in the layer.
\item \emph{Compatible coordinates.} Source and target nodes are typed
coordinate identities of quantity kind \texttt{DRIVE\_INTERVENTION} in
\textsc{alpha} units, and every node in the chain binds the same
registered test-point digest; a component whose topology nodes and whose
resistance measurement bind different test points is rejected.
\item \emph{Simple acyclic directed chain.} Adjacent components must link
exactly --- the target node of $C_i$ equals the source node of
$C_{i+1}$ --- and the full node path (first source, then each target in
order) must contain no repeated node. This excludes cycles, self-loops
(already rejected at the component level), branching, and re-visiting: the
topology is a simple directed path or nothing.
\item \emph{Active region.} Every component's declared region is
\texttt{ACTIVE}; the component type itself rejects any other region at
construction. Cutoff, saturated, non-ohmic, unstable, and unassessed
segments cannot enter a series computation at all.
\item \emph{Adequate linearity.} Every component's linearity coefficient
(an $R^2$ in $[0,1]$) is at or above the profile's declared minimum;
the threshold is a profile tolerance, not a constant.
\item \emph{One profile, with evidence.} Every component resistance is
validated against the single admitted calibration profile of the operation
context; all components bind one test point; each carries positive-lower-bound
uncertainty and a recorded sample size.
\item \emph{Complete ordered child provenance.} The emitted equivalent's
derivation parents are exactly the ordered component digests, in
role-stamped positions ($\texttt{component[0]}, \texttt{component[1]},
\ldots$); the result object re-verifies this binding at construction, and
any mismatch between claimed children and bound parents is a validation
error.
\end{enumerate}
\end{ruleenv}

Like topology, a component's declared region and linearity coefficient
are caller-asserted at this gate: the admission rules check enumeration
membership and threshold, not provenance. The gate does not require the
\texttt{ACTIVE} declaration to bind a replayable assessment
(Chapter~\ref{ch07:classification}) or the coefficient to derive from
the profile's fit receipt; a misdeclaring caller is undetectable at this
layer, and the limitation is declared rather than hidden.

When admitted, the arithmetic is the expected sum: the equivalent value is
$\sum_i r_i$; interval bounds sum lower-with-lower and upper-with-upper
(with the upper bound floored at the computed value); standard uncertainty
combines in quadrature, which silently encodes an independence assumption
that the covariance limitation of Section~\ref{ch09:sec:binding} makes
loud; the confidence level is the minimum across children; and the sample
size is the minimum across children. The sample-size convention is
conservative bookkeeping; the confidence convention is not conservative
--- the combined interval covers only when every child interval covers,
an event with probability at most $\min_i c_i$ (approximately
$\prod_i c_i$ under the declared zero-covariance approximation), so the
recorded level is an upper bound on joint coverage, and reading it as
attained coverage would overstate the evidence; replacing it with a
Bonferroni-style lower bound $\max\bigl(0,\,1-\sum_i(1-c_i)\bigr)$ is an
open revision item. The result is emitted with resistance role
\texttt{NETWORK\_EQUIVALENT}, which marks it as replayable-derived rather
than calibrated, and which downstream consumers (notably the prediction
gate) treat as a distinct role with its own replay obligation.

\subsection{The arithmetic, explicitly}

For admitted components with resistances $r_1, \ldots, r_n$, intervals
$[\ell_i, u_i]$, standard uncertainties $s_i$, confidence levels $c_i$,
and sample sizes $m_i$, the emitted equivalent is
\[
r_{\mathrm{series}} = \sum_{i=1}^{n} r_i, \qquad
\bigl[\;\textstyle\sum_i \ell_i,\; \max\bigl(r_{\mathrm{series}},
\sum_i u_i\bigr)\;\bigr], \qquad
s_{\mathrm{series}} = \Bigl(\sum_{i=1}^{n} s_i^2\Bigr)^{1/2},
\]
with confidence $\min_i c_i$ and sample size $\min_i m_i$. The
quadrature form for $s_{\mathrm{series}}$ is exact for uncorrelated
child errors and an approximation otherwise; since cross-input
covariance is unavailable in the current evidence model, the zero-covariance
approximation is declared on the binding rather than assumed silently
(Rule~\ref{ch09:rule:parallel-limits} applies to series results too).
If child errors were positively correlated, the true combined standard
uncertainty would exceed the quadrature value --- the reported
uncertainty would be optimistic --- and this direction of possible error
is precisely what the declared limitation warns a consumer to expect.

\begin{example}[Illustrative series computation]\label{ch09:ex:series}
All numbers illustrative. Two admitted components on one test point:
$r_1 = 2.0$ alpha-per-activation with interval $[1.8, 2.2]$, standard
uncertainty $0.1$, confidence $0.95$, sample size $40$; and $r_2 = 3.0$
with $[2.7, 3.3]$, standard uncertainty $0.15$, confidence $0.90$,
sample size $25$. Topology: $A \to B$, $B \to C$, distinct nodes, so the
chain is simple and acyclic. The equivalent is $r = 5.0$ with interval
$[4.5, 5.5]$, standard uncertainty $\sqrt{0.01 + 0.0225} \approx 0.18$,
confidence $0.90$ (the minimum), sample size $25$ (the minimum), role
\texttt{NETWORK\_EQUIVALENT}, derivation parents
$(\texttt{component[0]}, \texttt{component[1]})$ by digest, and the full
operation binding of Definition~\ref{ch09:def:binding}. If the same two
components were presented in the orders $(C_1, C_2)$ and $(C_2, C_1)$
with correspondingly re-declared node chains, the two results would be
distinct objects with distinct digests: child order is part of identity,
because the declared topology is directional.
\end{example}

\begin{remark}[What a series equivalent is not]\label{ch09:rem:series-not}
A series equivalent is not evidence that the two mechanisms are ``in
series'' in any causal sense; the chain topology is asserted by the
caller's component construction, not discovered. It is not a claim that
end-to-end response equals drive over summed resistance; that is the
content of a registered hypothesis, untested. And it is not transferable:
it binds one profile, one test point, and one replay frame, like every
other quantity in the candidate.
\end{remark}

\section{Parallel equivalence}\label{ch09:sec:parallel}

\subsection{Requirements as implemented}

Parallel equivalence shares the entire admission apparatus of
Rule~\ref{ch09:rule:series} --- typed components, calibrated or replayed
resistances, one unit, coordinate compatibility, active region, linearity
threshold, one profile, uncertainty and sample-size evidence, ordered child
binding, duplicate rejection --- and replaces the chain condition with an
endpoint condition:

\begin{ruleenv}[Parallel topology]\label{ch09:rule:parallel}
All components must share \emph{identical directed endpoints}: one
common source node and one common target node, as typed coordinate
identities on one registered test point. The set of (source, target) pairs
across components must have exactly one element. Components that merely
coexist, or that connect overlapping-but-distinct coordinates, do not form
a parallel topology and are rejected.
\end{ruleenv}

The admitted arithmetic is reciprocal aggregation:
\[
r_{\mathrm{parallel}} = \Bigl(\sum_{i=1}^{n} \frac{1}{r_i}\Bigr)^{-1},
\qquad
s_{\mathrm{parallel}} = \Biggl(\;\sum_{i=1}^{n}
\Bigl(\Bigl(\frac{r_{\mathrm{parallel}}}{r_i}\Bigr)^{2} s_i\Bigr)^{2}
\Biggr)^{1/2},
\]
with interval bounds computed by the same reciprocal form on the child
bounds, and the same minimum-confidence, minimum-sample-size,
role-\texttt{NETWORK\_EQUIVALENT}, factory-only, replay-bound conventions
as the series case. The weighting in $s_{\mathrm{parallel}}$ is the
first-order sensitivity of the reciprocal sum: a child whose resistance
is small relative to the equivalent dominates both the value and the
uncertainty, matching the electrical intuition that the lowest-resistance
branch dominates a parallel combination --- an intuition which, once
again, is inherited as arithmetic, not as validated mechanism.

\begin{example}[Illustrative parallel computation]\label{ch09:ex:parallel}
All numbers illustrative. The same two components as
Example~\ref{ch09:ex:series}, now declared with identical directed
endpoints $A \to B$ for both. The reciprocal sum is
$1/2.0 + 1/3.0 = 5/6$, so $r_{\mathrm{parallel}} = 1.2$ exactly in the
declared decimal precision, below both children, with interval bounds
$[\,1/(1/1.8 + 1/2.7),\; \max(1.2,\, 1/(1/2.2 + 1/3.3))\,] = [1.08, 1.32]$
and the propagated standard uncertainty computed from the squared-ratio
weights $(1.2/2.0)^2 = 0.36$ and $(1.2/3.0)^2 = 0.16$. Reading: under
the declared topology, \emph{if} the two components were genuine
alternative routes and \emph{if} the reciprocal law held empirically ---
the registered redundancy hypothesis, untested --- the aggregate
opposition would be $1.2$. Neither conditional is discharged by the
computation, and the result object says so through its binding.
\end{example}

\subsection{Co-present guardrails are not parallel paths}\label{ch09:sec:guardrails}

The electrical intuition behind reciprocal aggregation is that parallel
resistors offer \emph{alternative routes} for the same current: adding a
path can only decrease equivalent resistance, because flow that cannot
pass one way passes another. Imported uncritically into the safety domain,
this intuition produces a specific and consequential misreading, flagged
in Appendix~I and preserved in the source audit: that multiple co-present
guardrails --- several safety mechanisms bearing on the same behavior ---
constitute parallel resistances, so that adding a guardrail \emph{lowers}
total resistance and thereby \emph{weakens} the system. The arithmetic
would even look self-consistent, which is exactly what makes it dangerous.
The error is in the premise, not the algebra: co-presence is not
alternativeness.

Consider three scenarios, stated at the structural level the candidate can
actually represent.

\emph{Scenario 1: genuine alternatives.} Two mechanisms connect the same
source coordinate to the same target coordinate, and the response
suppressed along one route demonstrably appears along the other when the
first is ablated. This --- and only this --- is the situation the
reciprocal rule's intuition describes. Note that the italicized clause is
an experimental outcome, not a topological declaration: it is precisely
the registered redundancy hypothesis (ablation of genuine alternative
branches produces the specific response predicted by reciprocal
aggregation, above matched random-head controls), which is unproven. The
candidate will \emph{compute} the parallel equivalent for declared
same-endpoint components, but the computation's predictive meaning is
suspended pending that experiment.

\emph{Scenario 2: co-present guardrails on one pathway.} Three safety
mechanisms each independently gate the same route from drive to response
--- serially, as it were: each can block, none provides an alternative
route around another. Modeling these as parallel is not a poor
approximation; it is a category error with an inverted sign, since each
added guardrail here can only maintain or increase opposition to flow,
while the reciprocal rule asserts it decreases it. Structurally, the
candidate resists the error at the coordinate level: guardrails
instantiated at different layers or positions have different test-point or
coordinate identities and therefore fail the identical-endpoints gate.
Where a caller nonetheless declares them as same-endpoint components, the
type system cannot detect the misdeclaration --- topology is caller-asserted
--- and this is stated plainly as a limitation rather than hidden: the
parallel gate checks the declared topology's coherence, not its truth.

\emph{Scenario 3: correlated mechanisms.} Two mechanisms genuinely offer
alternative routes, but they share training provenance, circuitry, or
failure modes, so that conditions degrading one degrade the other. The
reciprocal arithmetic, and its quadrature uncertainty combination, assume
away exactly this correlation. The candidate cannot measure the
correlation (no cross-input covariance is available in the current
evidence model); what it can do is refuse to let the assumption pass
silently, which is the function of the mandatory covariance limitation
below.

\begin{ruleenv}[Declared limitations on every parallel result]\label{ch09:rule:parallel-limits}
Interaction, dependence, and covariance are not modeled by the parallel
operation and are declared, not waived: every network result's operation
binding carries the fixed covariance-limitation marker
(\texttt{CROSS\_INPUT\_COVARIANCE\_UNAVAILABLE\_ZERO\_COVARIANCE\_APPROXIMATION}),
and the binding validator rejects any binding in which the marker has been
altered or removed. A parallel equivalent thus cannot be excerpted without
excerpting its own disclaimer.
\end{ruleenv}

\section{Gain composition}\label{ch09:sec:gain}

\subsection{Requirements as implemented}

\begin{ruleenv}[Gain-composition admission]\label{ch09:rule:gain}
A gain composition over stages $S_1, \ldots, S_n$ ($n \ge 1$) is admitted
only when: each stage is an exact typed stage object with a unique stage
identity; each stage's gain is a \texttt{GAIN\_MAGNITUDE} in
\textsc{dimensionless} units with uncertainty and a recorded sample size,
validated against the context's single profile; each stage's source and
target coordinates are typed coordinate identities of \emph{identical}
quantity kind and unit binding one common test point (a dimensionless gain
may not map unlike coordinates); each stage's declared region is
\texttt{ACTIVE} (enforced at stage construction); each stage's linearity
coefficient meets the profile minimum; adjacent stages compose exactly ---
the target coordinate of $S_i$ equals the source coordinate of $S_{i+1}$
--- and the full coordinate path is simple and acyclic (no repeated
coordinate); and the emitted composite binds the ordered stage digests as
role-stamped derivation parents (\texttt{stage[0]}, \texttt{stage[1]},
\ldots), factory-only.
\end{ruleenv}

The admitted arithmetic is the product $\prod_i g_i$, with interval bounds
as products of child bounds, first-order product-rule propagation of
standard uncertainty (each child's contribution weighted by the product of
the others), minimum confidence, and minimum sample size. The composite is
emitted as a dimensionless gain magnitude named for its end-to-end
coordinate pair, with full provenance.

\subsection{The Fisher--Rao versus Jacobian distinction, and the missing bridge}\label{ch09:sec:gain-bridge}

The gain gate has a conceptual precondition that no runtime check can
discharge, preserved in the source audit as the gain-definition conflict.
The PCBA source defines stage gain $\beta$ by Fisher--Rao displacement
ratios on the open probability simplex: the magnitude of induced
displacement in an information-geometric metric, at small or large signal,
between cause and effect coordinates. The operationalization addendum
defines $\beta$ as a local Jacobian or projection derivative between cause
and effect test points: a ratio of infinitesimal changes in two activation
projections. These are distinct mathematical objects. One lives on a
statistical manifold and measures distributional displacement; the other
lives in activation space and measures a directional derivative. They have
different invariances (the Fisher--Rao quantity is
reparameterization-invariant on the simplex; the Jacobian quantity is
basis- and projection-dependent), different units discipline (both can be
rendered dimensionless, but by different normalizations), and no theorem
in the candidate's corpus identifies them --- no chain rule connects a
product of projection derivatives to a product of Fisher--Rao displacement
ratios except under assumptions (local diffeomorphism between coordinates,
metric compatibility) that nobody has stated, let alone verified.

The candidate's handling is deliberately unheroic: the two gains are
\emph{separately named profiles}. A gain magnitude enters the composition
gate carrying its method identity in provenance, and compositions are
computed within a method, under one profile. The missing bridge is
recorded as missing. What would supply it is one of two things: a theorem,
deriving one quantity from the other under stated conditions on the
coordinate maps; or the registered gain-composition experiment, which
tests empirically whether local derivative gains compose to end-to-end
gain in the active region within stated tolerance. Until one of these
exists, a composed gain is a product of numbers that were each,
individually, admitted under their own definition --- and the composite's
meaning is exactly as strong as the weakest link in that chain of
definitions, which is to say: structural.

\begin{conflict}[No licensed identification of gain profiles]\label{ch09:conf:gain}
Fisher--Rao displacement gain and local Jacobian/projection gain remain
distinct typed profiles in the candidate. No implemented operation
converts between them, no composition mixes them, and no result asserts
their equality. Any external text that multiplies a Fisher--Rao stage
figure by a Jacobian stage figure is performing an operation the candidate
refuses.
\end{conflict}

\section{What every aggregate binds}\label{ch09:sec:binding}

Every network result --- series, parallel, or gain --- embeds a provenance
structure and an operation binding whose fields are enumerated here,
because the binding \emph{is} the claim. A network equivalent, stripped of
its binding, is a number; with its binding, it is a replayable derivation
whose every premise is named.

\begin{definition}[Aggregate binding contents]\label{ch09:def:binding}
An admitted aggregate binds, at minimum:
\begin{enumerate}
\item \emph{Exact child order and identity} --- the ordered, role-stamped
digests of the components or stages, positionally contiguous from zero,
with duplicates rejected;
\item \emph{Topology} --- the declared chain or endpoint structure,
re-derivable from the typed children and re-verified on replay;
\item \emph{Profile identity} --- the calibration profile digest, binding
model, configuration, tokenizer, weights, code revision, runtime, dtype,
corpus, splits, seeds, alpha grid, full-scale method, fit and uncertainty
methods, tolerances, controls, drift and expiry, and review state;
\item \emph{Source and control evidence} --- the deduplicated source
artifact references of all children (with conflicting metadata for one
digest rejected), the profile's control references, and the profile's
control bindings, carried on the result's provenance;
\item \emph{Capability label} --- which gate was exercised
(\texttt{SERIES\_EQUIVALENCE}, \texttt{PARALLEL\_EQUIVALENCE},
\texttt{GAIN\_COMPOSITION});
\item \emph{Authority-manifest identity} --- the digest of the trust
authority's manifest under which the context was admitted;
\item \emph{Attestation identities} --- the sorted, deduplicated digests
of any capability attestations present in the context, \emph{possibly
empty}: structural mode carries none, and emptiness is a recorded fact,
not an omission;
\item \emph{Trust mode and authorization state} --- locked to each other
as described in Section~\ref{ch09:sec:philosophy}, both unauthenticated;
\item \emph{Historical replay time} --- the caller-supplied
\texttt{replay\_as\_of}, with the non-freshness semantics of
Chapter~\ref{ch07:classification}, plus the optional clock-receipt digest
when one was recorded;
\item \emph{Runtime, code, and dependency identities} --- the runtime
identifier, the code artifact reference, and the dependency-lock artifact
reference, each an exact typed reference with digest;
\item \emph{Covariance limitation} --- the fixed marker of
Rule~\ref{ch09:rule:parallel-limits}, mandatory on every binding, series
and gain included, because quadrature and product-rule propagation embed
the same zero-covariance approximation.
\end{enumerate}
\end{definition}

\begin{ruleenv}[Fail-closed topology]\label{ch09:rule:failclosed}
Duplicate children (by component identity, component digest, or underlying
resistance digest), repeated nodes in a series path, repeated coordinates
in a gain path, self-loops, non-contiguous parent positions, non-canonical
parent roles, unknown topology methods on replay, and any digest mismatch
between a claimed equivalent and its factory replay are validation errors.
There is no lenient mode, no warning-and-continue path, and no partial
aggregate: an aggregate that cannot bind everything binds nothing.
\end{ruleenv}

The replay obligation is the layer's
strongest structural property. A network-equivalent resistance presented
anywhere downstream --- as a child of a larger network, or as the
resistance input to the prediction gate --- is not consumed on the
strength of its role field. The consumer resolves its ordered parents from
the evidence registry, type-checks each as an exact component, re-runs the
declared topology method (series or parallel, and only those; an unknown
method identifier fails), and requires the replayed equivalent's digest to
equal the presented one. Nesting is therefore sound by induction: a tree
of aggregates validates only if every internal node replays exactly over
its leaves, and a single altered leaf invalidates every ancestor. The
uncertainty carried at each level remains a zero-covariance approximation
at that level --- the limitation marker propagates with the binding ---
so what the induction guarantees is integrity of derivation, not validity
of the composition law.

\section{Replay, and the prediction gate as consumer}\label{ch09:sec:replay}

The replay obligation sketched in Section~\ref{ch09:sec:philosophy} is
concrete enough to state as an algorithm, and doing so makes the layer's
inductive soundness inspectable.

\begin{algorithm}[Exact factory replay of a network equivalent]\label{ch09:alg:replay}
Input: an admitted operation context; a measurement value $R$ claimed to
be a network equivalent. Output: the replayed result, or a validation
error.
\begin{enumerate}
\item Require $R$ to be an exact measurement value with resistance role
\texttt{NETWORK\_EQUIVALENT}; validate $R$ against the context (unit
admissibility, profile binding, evidence resolution).
\item Read $R$'s derivation parents in order. For each position $i$,
require the parent role to be exactly $\texttt{component[}i\texttt{]}$;
resolve the parent digest through the evidence registry; require the
resolved object to be an exact resistance component. Any missing,
misrolled, misordered, or mistyped parent fails.
\item Dispatch on $R$'s provenance method identity: the series topology
method replays through the series factory, the parallel topology method
through the parallel factory; any other method identity fails as an
unknown topology.
\item Re-run the factory on the resolved children under the current
context. This re-executes every admission gate of
Rules~\ref{ch09:rule:series} and~\ref{ch09:rule:parallel} --- including
recursive replay of any child that is itself a network equivalent ---
and recomputes the arithmetic in the exact decimal context.
\item Require digest equality between the replayed equivalent and $R$.
Inequality --- from altered children, altered profile, altered
tolerances, or any drift whatsoever --- fails the replay.
\end{enumerate}
\end{algorithm}

Step 4 is where the induction closes: replay of a parent that is itself
an equivalent triggers the same algorithm one level down, so validating
the root of an aggregate tree validates every interior node against its
leaves. Step 5 is where integrity becomes binary: there is no
``approximately replayed'' outcome, and a single changed decimal
anywhere in the tree surfaces as root-level digest inequality.

The principal consumer of replayed equivalents is the prediction gate.
When a network-equivalent resistance is presented as the
$r_{\mathrm{eff}}$ of a candidate prediction $\hat y = a /
r_{\mathrm{eff}}$, the gate accepts the role --- calibrated-held-out and
network-equivalent are the only two roles it admits --- but only after
running Algorithm~\ref{ch09:alg:replay} to completion. The design
declines a tempting shortcut: it would have been cheaper to trust the
role tag and the stored digest. The gate instead re-derives, because the
role tag asserts a derivation history and the candidate's uniform rule
is that asserted history is checked history. A prediction computed from
a network equivalent therefore carries, transitively, the entire
admission record of every component in the tree --- and its own status
remains what Chapter~\ref{ch07:classification}'s boundary discipline
makes it: a structural research calculation whose empirical adequacy is
a registered, open question.

Table~\ref{ch09:tab:ops} summarizes the three operations side by side.

\begin{table}[htbp]
\centering
\small
\begin{tabular}{@{}llll@{}}
\toprule
 & Series & Parallel & Gain composition \\
\midrule
Children & $\ge 2$ components & $\ge 2$ components & $\ge 1$ stages \\
Child quantity & \texttt{EFFECTIVE\_RESISTANCE} & \texttt{EFFECTIVE\_RESISTANCE} & \texttt{GAIN\_MAGNITUDE} \\
Child unit & \texttt{ALPHA\_PER\_ACTIVATION} & \texttt{ALPHA\_PER\_ACTIVATION} & \texttt{DIMENSIONLESS} \\
Topology & simple directed chain & identical directed endpoints & simple coordinate chain \\
Region gate & \texttt{ACTIVE} only & \texttt{ACTIVE} only & \texttt{ACTIVE} only \\
Linearity gate & profile minimum & profile minimum & profile minimum \\
Value law & $\sum_i r_i$ & $(\sum_i r_i^{-1})^{-1}$ & $\prod_i g_i$ \\
Result role & \texttt{NETWORK\_EQUIVALENT} & \texttt{NETWORK\_EQUIVALENT} & (dimensionless gain) \\
Capability & \texttt{SERIES\_EQUIVALENCE} & \texttt{PARALLEL\_EQUIVALENCE} & \texttt{GAIN\_COMPOSITION} \\
Replayable & yes, by digest & yes, by digest & factory-only \\
\bottomrule
\end{tabular}
\caption{The three structural network operations. All three share the
one-profile, one-test-point, uncertainty-and-sample-size, ordered-child-binding,
duplicate-rejection, and covariance-limitation requirements.}
\label{ch09:tab:ops}
\end{table}

The gain column's \emph{factory-only} entry is a declared asymmetry, not
an oversight: composed gains are factory-bound at construction, but the
replay algorithm dispatches only the series and parallel resistance
methods, so gain compositions have no digest-replay verifier in this
candidate; closing that gap is an open item.

\section{The Ed25519 fixture-key mode}\label{ch09:sec:ed25519}

The candidate's second trust mode, \texttt{INTEGRITY\_CONFIGURATION},
admits optional capability attestations: Ed25519 signatures over
canonical attestation payloads, verified against a capability authority's
manifest, checked for capability match, profile binding, validity window
against the replay time, and evidence presence, and replayed like
everything else. The keys involved are \emph{fixture keys}: deterministic
in-repository test keys, present so that the signature-verification code
path is real, exercised, and bounded by tests, rather than a stub awaiting
hardware.

What the mode is for: integrity and replay discipline. A signed
attestation binds a capability grant to a manifest digest, a profile
digest, and a validity window, so that tests can verify the candidate
rejects expired attestations, wrong-capability attestations,
wrong-profile attestations, tampered payloads, and unknown signers ---
the whole negative space of a capability system --- with real
cryptographic verification rather than mocks.

The verification path, when the mode is engaged, checks in sequence:
that an attestation for the requested capability is present in the
context (absence is a capability error, not a silent downgrade); that
the signature verifies over the attestation's canonical bytes against
the key identified in the authority manifest; that the attestation binds
the context's exact profile digest; that the replay time falls inside
the attestation's validity window; and that the attestation's evidence
reference resolves in the registry. Each check fails closed, and the
attestation digests --- sorted, deduplicated --- enter the operation
binding of every result computed under the context, so that a replayed
result re-verifies the same attestations it was computed under. In
structural mode, by contrast, the attestation tuple is empty, the
integrity requirement is simply not asserted, and the binding records
the emptiness. The two modes thus differ observably, in their recorded
bindings, while agreeing on the only point that matters for authority:
neither is authenticated.

What the mode is not, stated with the bluntness the source material uses:

\begin{ruleenv}[Fixture keys confer no authority]\label{ch09:rule:fixture}
Verification of an Ed25519 fixture-key attestation establishes only that
a payload was signed by a key present in the repository's test
configuration and that the structural checks pass. It is never
independent reviewer authority, never legal authority, never product
authority, and never Evidence-Packet authority. The authorization state
of an integrity-mode operation remains
\texttt{UNAUTHENTICATED\_INTEGRITY} --- unauthenticated by type --- and
authenticated Evidence-Packet eligibility remains constant false in this
candidate. Deterministic in-repository keys and caller-supplied replay
time can never satisfy the eligibility preconditions (independently
admitted authority manifest, protected key lifecycle and revocation,
trusted current-time receipt, materialized dependency identities,
authenticated instrument evidence, reviewed packet policy). No
authenticated capability path is implemented: there is no code in the
candidate that, given better keys, would begin asserting authority ---
the assertion itself is absent from the type system.
\end{ruleenv}

The design rationale mirrors the completeness-versus-eligibility split
that runs through the whole candidate. It would have been easy, and
misleading, to ship the network layer with a configuration flag named
something like ``trusted mode.'' The implemented alternative is a pair of
modes that differ in \emph{which integrity checks run}, not in \emph{what
authority results carry}: both are unauthenticated, and the enumeration
contains nothing else. A future authenticated path, if ever built, would
be a new admission --- new types, new review, new gates --- not a
configuration change to this one.

\section{Chapter boundary}\label{ch09:sec:boundary}

\begin{boundary}[What the network layer establishes]\label{ch09:bound:networks}
The network layer establishes that series, parallel, and gain arithmetic
can be performed under exhaustive structural gating with complete,
replayable provenance, and that every result carries its own limitations
--- unauthenticated status, zero-covariance approximation, profile
relativity, caller-asserted topology --- inseparably. It establishes
nothing about whether steered activation pathways compose like circuits.
The Guide's network rules remain, in this candidate, preregistered
predictions with named falsifiers; the redundancy and gain-composition
experiments have not been run; and every equivalent resistance and
composed gain in existence under this candidate is a structural research
calculation, Ring~4, not licensed, held.
\end{boundary}
